% GENERATED FILE. Edit canonical lessons, then run npm run generate:books.
% canonical-source-hash: fnv1a64:5aab75ad67a7cebb
% canonical-lessons: ML-C461-plarrphom, ML-C461-kantaktar, ML-C461-yatrakkuli, ML-C461-kyu, ML-C461-sancarikkuka

\chapter{At the Platform}
\label{ch:ml-at-the-platform}
\hlchaptermodality{461}

\begin{chapteropening}
\textbf{By the end of this chapter:} \emph{I can say a platform, a bus conductor, a fare, a queue and to travel.}

Complete the last lesson of chapter 461: I can say a platform, a bus conductor, a fare, a queue and to travel.
\end{chapteropening}

\section[\texorpdfstring{\ml{പ്ലാറ്റ്ഫോം} (plāṟṟphōṁ)}{പ്ലാറ്റ്ഫോം (plāṟṟphōṁ)}]{\ml{പ്ലാറ്റ്ഫോം} (plāṟṟphōṁ) — a platform}

\label{lesson:ML-C461-plarrphom}



\begin{quote}
Before the new one: say the Malayalam for a principal, then the Malayalam for tuition.
\end{quote}

\subsection*{You'll want to know: \ml{പ്ലാറ്റ്ഫോം}}
\textbf{\ml{പ്ലാറ്റ്ഫോം}} — \emph{plāṟṟphōṁ} — \textquotedblleft{}a platform\textquotedblright{}.

\begin{quote}
\ml{തീവണ്ടി} \ml{മൂന്നാം} \ml{പ്ലാറ്റ്ഫോമിൽ} \ml{എത്തി}.
\end{quote}

\begin{quote}
\emph{tīvaṇṭi mūnnāṁ plāṟṟphōmil etti.}
\end{quote}

\begin{quote}
The train arrived at platform three.
\end{quote}

Platforms are counted with the ordinal: \ml{ഒന്നാം} \ml{പ്ലാറ്റ്ഫോം}, platform one; \ml{രണ്ടാം} \ml{പ്ലാറ്റ്ഫോം}, platform two.

\begin{cousinweb}
\textbf{\ml{പ്ലാറ്റ്ഫോം}} (\emph{plāṟṟphōṁ}) was borrowed from English \emph{platform}, the raised floor beside a railway line.
\end{cousinweb}

\begin{grammarlens}[title={What you've built}]
One more word for getting around.
\end{grammarlens}

\subsection*{Your turn}
\begin{itemize}
\raggedright
  \item \emph{Say it:} \emph{plāṟṟphōṁ}
  \item \emph{Say it:} \emph{plāṟṟphōṁ}, once more
  \item \emph{Say it:} \emph{tīvaṇṭi mūnnāṁ plāṟṟphōmil etti.} — The train arrived at platform three.
  \item \emph{Recall:} say the Malayalam for a principal, then the Malayalam for tuition, then say \emph{plāṟṟphōṁ} again
\end{itemize}

\begin{tcolorbox}[breakable,colback=teal!4,colframe=teal!35!black,title={Before you move on}]
What does \ml{പ്ലാറ്റ്ഫോം} mean? (\textquotedblleft{}A platform\textquotedblright{}.) Say it once more.
\end{tcolorbox}
\section[\texorpdfstring{\ml{കണ്ടക്ടർ} (kaṇṭakṭar)}{കണ്ടക്ടർ (kaṇṭakṭar)}]{\ml{കണ്ടക്ടർ} (kaṇṭakṭar) — a bus conductor}

\label{lesson:ML-C461-kantaktar}



\begin{quote}
Before the new one: say the Malayalam for tuition, then the Malayalam for a platform.
\end{quote}

\subsection*{You'll want to know: \ml{കണ്ടക്ടർ}}
\textbf{\ml{കണ്ടക്ടർ}} — \emph{kaṇṭakṭar} — \textquotedblleft{}a bus conductor\textquotedblright{}.

\begin{quote}
\ml{കണ്ടക്ടർ} \ml{എനിക്ക്} \ml{ടിക്കറ്റ്} \ml{തന്നു}.
\end{quote}

\begin{quote}
\emph{kaṇṭakṭar enikkŭ ṭikkaṟṟŭ tannu.}
\end{quote}

\begin{quote}
The conductor gave me a ticket.
\end{quote}

On a bus the conductor sells the tickets and gives the change; the one who drives is a separate person.

\begin{cousinweb}
\textbf{\ml{കണ്ടക്ടർ}} (\emph{kaṇṭakṭar}) was borrowed from English \emph{conductor}, the person who takes the fares on a bus or tram.
\end{cousinweb}

\begin{grammarlens}[title={What you've built}]
One more word for getting around.
\end{grammarlens}

\subsection*{Your turn}
\begin{itemize}
\raggedright
  \item \emph{Say it:} \emph{kaṇṭakṭar}
  \item \emph{Say it:} \emph{kaṇṭakṭar}, once more
  \item \emph{Say it:} \emph{kaṇṭakṭar enikkŭ ṭikkaṟṟŭ tannu.} — The conductor gave me a ticket.
  \item \emph{Recall:} say the Malayalam for tuition, then the Malayalam for a platform, then say \emph{kaṇṭakṭar} again
\end{itemize}

\begin{tcolorbox}[breakable,colback=teal!4,colframe=teal!35!black,title={Before you move on}]
What does \ml{കണ്ടക്ടർ} mean? (\textquotedblleft{}A bus conductor\textquotedblright{}.) Say it once more.
\end{tcolorbox}
\section[\texorpdfstring{\ml{യാത്രാക്കൂലി} (yātrākkūli)}{യാത്രാക്കൂലി (yātrākkūli)}]{\ml{യാത്രാക്കൂലി} (yātrākkūli) — a fare}

\label{lesson:ML-C461-yatrakkuli}



\begin{quote}
Before the new one: say the Malayalam for a platform, then the Malayalam for a bus conductor.
\end{quote}

\subsection*{You'll want to know: \ml{യാത്രാക്കൂലി}}
\textbf{\ml{യാത്രാക്കൂലി}} — \emph{yātrākkūli} — \textquotedblleft{}a fare\textquotedblright{}.

\begin{quote}
\ml{ബസിലെ} \ml{യാത്രാക്കൂലി} \ml{ഇപ്പോൾ} \ml{കൂടുതലാണ്}.
\end{quote}

\begin{quote}
\emph{basile yātrākkūli ippōḷ kūṭutalāṇŭ.}
\end{quote}

\begin{quote}
The bus fare is higher now.
\end{quote}

Built from \ml{യാത്ര}, a journey, and \ml{കൂലി}, wages: the price of a journey. Before a joined word \ml{യാത്ര} takes a long -\ml{ആ}.

\begin{cousinweb}
\textbf{\ml{യാത്രാക്കൂലി}} (\emph{yātrākkūli}) is a Malayalam compound built from \ml{യാത്ര}, a journey, borrowed from Sanskrit \emph{yātrā} (CDIAL 10456), and the native word \ml{കൂലി}, hire or wages (DEDR 1905).
\end{cousinweb}

\begin{grammarlens}[title={What you've built}]
One more word for getting around.
\end{grammarlens}

\subsection*{Your turn}
\begin{itemize}
\raggedright
  \item \emph{Say it:} \emph{yātrākkūli}
  \item \emph{Say it:} \emph{yātrākkūli}, once more
  \item \emph{Say it:} \emph{basile yātrākkūli ippōḷ kūṭutalāṇŭ.} — The bus fare is higher now.
  \item \emph{Recall:} say the Malayalam for a platform, then the Malayalam for a bus conductor, then say \emph{yātrākkūli} again
\end{itemize}

\begin{tcolorbox}[breakable,colback=teal!4,colframe=teal!35!black,title={Before you move on}]
What does \ml{യാത്രാക്കൂലി} mean? (\textquotedblleft{}A fare\textquotedblright{}.) Say it once more.
\end{tcolorbox}
\section[\texorpdfstring{\ml{ക്യൂ} (kyū)}{ക്യൂ (kyū)}]{\ml{ക്യൂ} (kyū) — a queue}

\label{lesson:ML-C461-kyu}



\begin{quote}
Before the new one: say the Malayalam for a bus conductor, then the Malayalam for a fare.
\end{quote}

\subsection*{You'll want to know: \ml{ക്യൂ}}
\textbf{\ml{ക്യൂ}} — \emph{kyū} — \textquotedblleft{}a queue\textquotedblright{}.

\begin{quote}
\ml{ടിക്കറ്റിനായി} \ml{വലിയ} \ml{ക്യൂ} \ml{ഉണ്ടായിരുന്നു}.
\end{quote}

\begin{quote}
\emph{ṭikkaṟṟināyi valiya kyū uṇṭāyirunnu.}
\end{quote}

\begin{quote}
There was a long queue for tickets.
\end{quote}

You stand in a queue: \ml{ക്യൂവിൽ} \ml{നിൽക്കുക}.

\begin{cousinweb}
\textbf{\ml{ക്യൂ}} (\emph{kyū}) was borrowed from English \emph{queue}, a line of people waiting their turn.
\end{cousinweb}

\begin{grammarlens}[title={What you've built}]
One more word for getting around.
\end{grammarlens}

\subsection*{Your turn}
\begin{itemize}
\raggedright
  \item \emph{Say it:} \emph{kyū}
  \item \emph{Say it:} \emph{kyū}, once more
  \item \emph{Say it:} \emph{ṭikkaṟṟināyi valiya kyū uṇṭāyirunnu.} — There was a long queue for tickets.
  \item \emph{Recall:} say the Malayalam for a bus conductor, then the Malayalam for a fare, then say \emph{kyū} again
\end{itemize}

\begin{tcolorbox}[breakable,colback=teal!4,colframe=teal!35!black,title={Before you move on}]
What does \ml{ക്യൂ} mean? (\textquotedblleft{}A queue\textquotedblright{}.) Say it once more.
\end{tcolorbox}
\section[\texorpdfstring{\ml{സഞ്ചരിക്കുക} (sañcarikkuka)}{സഞ്ചരിക്കുക (sañcarikkuka)}]{\ml{സഞ്ചരിക്കുക} (sañcarikkuka) — to travel; to move about}

\label{lesson:ML-C461-sancarikkuka}



\begin{quote}
Before the new one: say the Malayalam for a fare, then the Malayalam for a queue.
\end{quote}

\subsection*{You'll want to know: \ml{സഞ്ചരിക്കുക}}
\textbf{\ml{സഞ്ചരിക്കുക}} — \emph{sañcarikkuka} — \textquotedblleft{}to travel; to move about\textquotedblright{}.

\begin{quote}
\ml{അവർ} \ml{തീവണ്ടിയിൽ} \ml{പല} \ml{നഗരങ്ങളിലും} \ml{സഞ്ചരിച്ചു}.
\end{quote}

\begin{quote}
\emph{avar tīvaṇṭiyil pala nagaraṅṅaḷiluṁ sañcariccu.}
\end{quote}

\begin{quote}
They travelled to many cities by train.
\end{quote}

Past: \ml{സഞ്ചരിച്ചു}, \emph{sañcariccu}.

A more formal word than \ml{യാത്ര} \ml{ചെയ്യുക}. The noun is \ml{സഞ്ചാരം}, travel, and one who travels is a \ml{സഞ്ചാരി}.

\begin{cousinweb}
\textbf{\ml{സഞ്ചരിക്കുക}} (\emph{sañcarikkuka}) is built from Sanskrit \emph{saṁcara}, \textquotedblleft{}going about, moving\textquotedblright{} (Monier-Williams), with the Malayalam verb ending -\ml{ിക്കുക}, the ending many Sanskrit nouns take to become verbs.
\end{cousinweb}

\begin{grammarlens}[title={What you've built}]
One more word for getting around.
\end{grammarlens}

\subsection*{Your turn}
\begin{itemize}
\raggedright
  \item \emph{Say it:} \emph{sañcarikkuka}
  \item \emph{Say it:} \emph{sañcarikkuka}, once more
  \item \emph{Say it:} \emph{avar tīvaṇṭiyil pala nagaraṅṅaḷiluṁ sañcariccu.} — They travelled to many cities by train.
  \item \emph{Recall:} say the Malayalam for a fare, then the Malayalam for a queue, then say \emph{sañcarikkuka} again
\end{itemize}

\begin{tcolorbox}[breakable,colback=teal!4,colframe=teal!35!black,title={Before you move on}]
What does \ml{സഞ്ചരിക്കുക} mean? (\textquotedblleft{}To travel; to move about\textquotedblright{}.) Say it once more.
\end{tcolorbox}
